\documentclass[12pt]{article}
\usepackage{amsmath, amssymb, amsthm, xcolor, mdframed, array}
\usepackage{geometry}
\usepackage{titlesec}
\usepackage{parskip}
\usepackage{microtype}
\usepackage{enumitem}

\geometry{margin=1.3in, top=1.2in, bottom=1.4in}

%% ---------------------------------------------------------------
%%  Colors
%% ---------------------------------------------------------------
\definecolor{exampleblue}{RGB}{30, 90, 180}
\definecolor{examplebg}{RGB}{235, 243, 255}
\definecolor{motivationgreen}{RGB}{20, 110, 60}
\definecolor{motivationbg}{RGB}{235, 248, 240}
\definecolor{headergray}{RGB}{60, 60, 60}
\definecolor{stepcolor}{RGB}{120, 0, 0}

%% ---------------------------------------------------------------
%%  Section / heading styles
%% ---------------------------------------------------------------
\titleformat{\section}
  {\large\bfseries\color{headergray}}
  {}{}{}[\vspace{2pt}\titlerule\vspace{4pt}]

\titleformat{\subsection}
  {\normalsize\bfseries\color{stepcolor}}
  {}{}{}

\titlespacing{\section}{0pt}{24pt}{10pt}
\titlespacing{\subsection}{0pt}{18pt}{6pt}

%% ---------------------------------------------------------------
%%  Theorem style
%% ---------------------------------------------------------------
\newtheoremstyle{mythmstyle}
  {14pt}{14pt}
  {\itshape}{}
  {\bfseries}{.}
  { }{}
\theoremstyle{mythmstyle}
\newtheorem{theorem}{Theorem}

%% ---------------------------------------------------------------
%%  Colored environments
%% ---------------------------------------------------------------
\newenvironment{example}{%
  \vspace{8pt}
  \begin{mdframed}[backgroundcolor=examplebg, linecolor=exampleblue,
    linewidth=1.4pt, innerleftmargin=14pt, innerrightmargin=14pt,
    innertopmargin=12pt, innerbottommargin=12pt,
    skipabove=12pt, skipbelow=12pt]
  \noindent\textcolor{exampleblue}{\textbf{\small CONCRETE EXAMPLE} \normalsize($V = \mathbb{R}^2$).}\par\vspace{6pt}
}{%
  \end{mdframed}
  \vspace{6pt}
}

\newenvironment{motivation}{%
  \vspace{6pt}
  \begin{mdframed}[backgroundcolor=motivationbg, linecolor=motivationgreen,
    linewidth=1.4pt, innerleftmargin=14pt, innerrightmargin=14pt,
    innertopmargin=12pt, innerbottommargin=12pt,
    skipabove=12pt, skipbelow=12pt]
  \noindent\textcolor{motivationgreen}{\textbf{\small MOTIVATION \& SCOPE}}\par\vspace{6pt}
}{%
  \end{mdframed}
  \vspace{6pt}
}

%% ---------------------------------------------------------------
%%  Paragraph spacing
%% ---------------------------------------------------------------
\setlength{\parskip}{9pt}
\setlength{\parindent}{0pt}

\begin{document}

%% ---------------------------------------------------------------
%%  Title block
%% ---------------------------------------------------------------
\begin{center}
  {\LARGE\bfseries\color{headergray} Two-Sided Ideals of $\mathcal{L}(V)$}

  \vspace{14pt}

  {\large Liam Abrams}

  \vspace{4pt}

  {\normalsize \today}

  \vspace{10pt}

  {\normalsize\itshape
    \textit{Linear Algebra Done Right, 4th Ed.}, Sheldon Axler\\[2pt]
    Chapter 3 \ $\cdot$ \ Problem 17}
\end{center}

\vspace{20pt}
\hrule height 0.8pt
\vspace{20pt}

%% ---------------------------------------------------------------
%%  Motivation
%% ---------------------------------------------------------------
\begin{motivation}
The goal of this writeup is to prove that the only two-sided ideals of
$\mathcal{L}(V)$ are $\{0\}$ and $\mathcal{L}(V)$, using \emph{only} the
theory developed through Chapter~3 of Axler --- that is, the basic properties
of vector spaces, subspaces, bases, and linear maps. We deliberately avoid
ring theory and any machinery from later in the book. The key tools are:
\begin{itemize}[leftmargin=1.4em, itemsep=3pt]
  \item A linear map is freely and uniquely determined by its values on a basis.
  \item A subspace is closed under addition.
\end{itemize}
\end{motivation}

%% ---------------------------------------------------------------
%%  Key definition
%% ---------------------------------------------------------------
\section{Key Definition}

A subspace $\mathcal{E} \subseteq \mathcal{L}(V)$ is called a
\emph{two-sided ideal} if for every $E \in \mathcal{E}$ and every
$T \in \mathcal{L}(V)$,
\[
    TE \in \mathcal{E} \qquad \text{and} \qquad ET \in \mathcal{E}.
\]

\vspace{4pt}

In words: you can multiply any element of $\mathcal{E}$ on the \emph{left or
right} by any operator in $\mathcal{L}(V)$, and the result always stays
inside $\mathcal{E}$.

%% ---------------------------------------------------------------
%%  Theorem
%% ---------------------------------------------------------------
\section{Main Theorem}

\begin{theorem}
Suppose $V$ is finite-dimensional. The only two-sided ideals of
$\mathcal{L}(V)$ are $\{0\}$ and $\mathcal{L}(V)$.
\end{theorem}

%% ---------------------------------------------------------------
%%  Proof
%% ---------------------------------------------------------------
\section{Proof}

Let $\mathcal{E}$ be a two-sided ideal of $\mathcal{L}(V)$. It is trivial to show that ${0}$ is a two-sided ideal, so assume $\mathcal{E}$ is nonzero.
We will show $\mathcal{E} = \mathcal{L}(V)$ by showing that \emph{every}
operator in $\mathcal{L}(V)$ must lie in $\mathcal{E}$.

%% --- Step 1 ---
\subsection{Step 1 \quad Find a vector that $E$ does not kill}

Since $\mathcal{E} \neq \{0\}$, there exists a nonzero operator
$E \in \mathcal{E}$. Since $E$ is not the zero map, there must exist some
vector $u \in V$ such that $Eu \neq 0$. This is the one foothold we have
inside $\mathcal{E}$: a single nonzero operator $E$, and a single vector
$u$ that $E$ moves somewhere nontrivial.

\vspace{6pt}

Extend $u$ to a basis $\,u,\, f_2,\, \ldots,\, f_n\,$ of $V$.\\
Extend $Eu$ to a basis $\,Eu,\, g_2,\, \ldots,\, g_n\,$ of $V$.

\begin{example}
Take $V = \mathbb{R}^2$ and suppose $\mathcal{E}$ contains the nonzero
operator
\[
    E = \begin{pmatrix} 1 & 0 \\ 0 & 0 \end{pmatrix}.
\]

Try $u = e_1 = \begin{pmatrix}1\\0\end{pmatrix}$. Then
$Eu = \begin{pmatrix}1\\0\end{pmatrix} \neq 0$. \checkmark

\vspace{6pt}

The basis built around $u$ is the standard basis:
\[
  u = e_1 = \begin{pmatrix}1\\0\end{pmatrix}, \qquad
  f_2 = e_2 = \begin{pmatrix}0\\1\end{pmatrix}.
\]
Since $Eu = e_1$, the basis built around $Eu$ is also $\,Eu = e_1,\;
g_2 = e_2$.
\end{example}

%% --- Step 2 ---
\subsection{Step 2 \quad Build a ``controlled'' operator using the ideal}

The goal of this step is to show that for \emph{any} basis vector $v_i$
and any target vector $w$, there exists an operator in $\mathcal{E}$ that
sends $v_i \mapsto w$ and kills all other basis vectors. We do this by
sandwiching $E$ between two hand-crafted operators.

\vspace{6pt}

Let $v_1, \ldots, v_n$ be any basis of $V$, let $F \in \mathcal{L}(V)$
be any operator, and fix an index $i$. We want to produce an operator
$A_i \in \mathcal{E}$ satisfying
\[
    A_i(v_i) = Fv_i \qquad \text{and} \qquad A_i(v_j) = 0
    \;\text{ for all }\; j \neq i.
\]

\vspace{4pt}

Construct two operators $R$ and $S$ as follows.

\vspace{10pt}

\noindent\textbf{Define $R$ \ (the ``router in''):}
\[
    R(v_i) = u, \qquad R(v_j) = 0 \;\text{ for all }\; j \neq i.
\]
$R$ is defined freely on the basis $v_1, \ldots, v_n$, so it is a valid
linear map. Its job is to funnel the chosen input $v_i$ into $u$ --- the
one vector we know $E$ handles nontrivially.

\vspace{10pt}

\noindent\textbf{Define $S$ \ (the ``router out''):}
\[
    S(Eu) = Fv_i, \qquad S(g_k) = 0 \;\text{ for all }\; k \geq 2.
\]
$S$ is defined freely on the basis $Eu, g_2, \ldots, g_n$, so it is a
valid linear map. Its job is to send $Eu$ (the output of $E$) to whatever
target we want, namely $Fv_i$.

\vspace{12pt}

\noindent Now set $A_i = SER$ and verify that $A_i \in \mathcal{E}$
using the ideal property twice:
\begin{align*}
    E \in \mathcal{E}
    &\implies ER \in \mathcal{E}
    &&\text{($\mathcal{E}$ closed under right-multiplication by $R$)},\\[6pt]
    ER \in \mathcal{E}
    &\implies S(ER) = A_i \in \mathcal{E}
    &&\text{($\mathcal{E}$ closed under left-multiplication by $S$)}.
\end{align*}

\vspace{4pt}

Finally, trace each basis vector through the chain
$\;v_j \xrightarrow{R} \cdot \xrightarrow{E} \cdot \xrightarrow{S} \cdot\;$
to confirm $A_i$ behaves as claimed:
\begin{align*}
    A_i(v_i) &= S\!\left(E\!\left(R(v_i)\right)\right)
               = S(E(u)) = S(Eu) = Fv_i, \\[6pt]
    A_i(v_j) &= S\!\left(E\!\left(R(v_j)\right)\right)
               = S(E(0)) = S(0) = 0
               \qquad \text{for all } j \neq i.
\end{align*}

\begin{example}
Let the target basis be the standard basis
$v_1 = e_1,\; v_2 = e_2$, and let
\[
  F = \begin{pmatrix}3&7\\2&5\end{pmatrix}
\]
be any operator we wish to recover.

\vspace{10pt}

\noindent\textbf{Building $A_1$} \quad (handles $v_1 = e_1$;
target $Fe_1 = \begin{pmatrix}3\\2\end{pmatrix}$, kills $e_2$)

\vspace{6pt}

\noindent\emph{Define $R_1$:} \quad
$R_1(e_1) = u = e_1,\quad R_1(e_2) = 0.$ \quad
As a matrix: $R_1 = \begin{pmatrix}1&0\\0&0\end{pmatrix}$.

\vspace{4pt}

\noindent\emph{Define $S_1$:} \quad
$S_1(Eu) = S_1(e_1) = \begin{pmatrix}3\\2\end{pmatrix},\quad S_1(e_2) = 0.$ \quad
As a matrix: $S_1 = \begin{pmatrix}3&0\\2&0\end{pmatrix}$.

\vspace{4pt}

\noindent\emph{Compute $A_1 = S_1 E R_1$:}
\[
    A_1 =
    \begin{pmatrix}3&0\\2&0\end{pmatrix}
    \begin{pmatrix}1&0\\0&0\end{pmatrix}
    \begin{pmatrix}1&0\\0&0\end{pmatrix}
    =
    \begin{pmatrix}3&0\\2&0\end{pmatrix}.
\]
\noindent Check: \;
$A_1(e_1) = \begin{pmatrix}3\\2\end{pmatrix} = Fe_1$ \checkmark \qquad
$A_1(e_2) = \begin{pmatrix}0\\0\end{pmatrix}$ \checkmark

\vspace{14pt}

\noindent\textbf{Building $A_2$} \quad (handles $v_2 = e_2$;
target $Fe_2 = \begin{pmatrix}7\\5\end{pmatrix}$, kills $e_1$)

\vspace{6pt}

\noindent\emph{Define $R_2$:} \quad
$R_2(e_1) = 0,\quad R_2(e_2) = u = e_1.$ \quad
As a matrix: $R_2 = \begin{pmatrix}0&1\\0&0\end{pmatrix}$.

\vspace{4pt}

\noindent\emph{Define $S_2$:} \quad
$S_2(Eu) = S_2(e_1) = \begin{pmatrix}7\\5\end{pmatrix},\quad S_2(e_2) = 0.$ \quad
As a matrix: $S_2 = \begin{pmatrix}7&0\\5&0\end{pmatrix}$.

\vspace{4pt}

\noindent\emph{Compute $A_2 = S_2 E R_2$:}
\[
    A_2 =
    \begin{pmatrix}7&0\\5&0\end{pmatrix}
    \begin{pmatrix}1&0\\0&0\end{pmatrix}
    \begin{pmatrix}0&1\\0&0\end{pmatrix}
    =
    \begin{pmatrix}0&7\\0&5\end{pmatrix}.
\]
\noindent Check: \;
$A_2(e_1) = \begin{pmatrix}0\\0\end{pmatrix}$ \checkmark \qquad
$A_2(e_2) = \begin{pmatrix}7\\5\end{pmatrix} = Fe_2$ \checkmark
\end{example}

%% --- Step 3 ---
\subsection{Step 3 \quad Reconstruct any operator $F$ as a sum of the $A_i$}

Applying Step~2 for each index $i = 1, \ldots, n$ produces operators
$A_1, \ldots, A_n \in \mathcal{E}$, each handling exactly one basis vector
and killing the rest. We claim these pieces add up to $F$ exactly.

\vspace{6pt}

Evaluate $A_1 + \cdots + A_n$ on an arbitrary basis vector $v_i$:
\[
    (A_1 + \cdots + A_n)(v_i)
    \;=\;
    \underbrace{A_1(v_i)}_{=\,0}
    + \cdots +
    \underbrace{A_{i-1}(v_i)}_{=\,0}
    +
    \underbrace{A_i(v_i)}_{=\,Fv_i}
    +
    \underbrace{A_{i+1}(v_i)}_{=\,0}
    + \cdots +
    \underbrace{A_n(v_i)}_{=\,0}
    \;=\; Fv_i.
\]

\vspace{4pt}

Since $F$ and $A_1 + \cdots + A_n$ agree on every basis vector, they are
the same linear map:
\[
    \boxed{F \;=\; A_1 \;+\; A_2 \;+\; \cdots \;+\; A_n.}
\]

Since $\mathcal{E}$ is a subspace, it is closed under addition. Each
$A_i \in \mathcal{E}$, so their sum $F \in \mathcal{E}$.

\begin{example}
\[
    A_1 + A_2
    \;=\;
    \begin{pmatrix}3&0\\2&0\end{pmatrix}
    +
    \begin{pmatrix}0&7\\0&5\end{pmatrix}
    \;=\;
    \begin{pmatrix}3&7\\2&5\end{pmatrix}
    \;=\; F. \qquad \checkmark
\]

\vspace{4pt}

So $F \in \mathcal{E}$, even though we only assumed the single matrix
$E = \begin{pmatrix}1&0\\0&0\end{pmatrix}$ was in $\mathcal{E}$.
\end{example}

%% ---------------------------------------------------------------
%%  Conclusion
%% ---------------------------------------------------------------
\section{Conclusion}

Since $F \in \mathcal{L}(V)$ was completely arbitrary, every operator
lies in $\mathcal{E}$, giving $\mathcal{L}(V) \subseteq \mathcal{E}$.
Since $\mathcal{E}$ is by definition a subspace of $\mathcal{L}(V)$, we
also have $\mathcal{E} \subseteq \mathcal{L}(V)$. Therefore
\[
    \mathcal{E} = \mathcal{L}(V). \qquad \blacksquare
\]

\end{document}
